\section{Setup: estimand, estimator, and where its error comes from}
\label{sec:setup}
\textbf{Estimand and model-defined truth.}
A graph $G$ on $n$ units has adjacency $A$, degrees $d_i$ and row-normalised $W=D^{-1}A$. An assignment $z\in\{0,1\}^n$ induces the
neighbour-treated fraction $f_i(z)=(Wz)_i$. The potential outcome $Y_i(z)$ may depend on the whole of $z$. The global average treatment effect is
$\tau_G=\frac1n\sum_i\,[Y_i(\mathbf 1)-Y_i(\mathbf 0)]$. Because we \emph{define} every outcome function, $\tau_G$ is computed exactly for each
simulated world (graph draw, baselines and noise); it is never estimated. All errors below are $\hat\tau-\tau_G$.

\textbf{Outcome mechanisms.}
The default is one-hop spillover $Y_i(z)=b_i+\tau z_i+\gamma\,h(f_i(z))+\varepsilon_i$ with direct effect $\tau=1$, strength $\gamma$, shape $h$ ($h(0)=0$, $h(1)=1$, so $\tau_G=\tau+\gamma$),
a baseline $b_i$ (graph-smooth, log-degree and i.i.d.\ parts) and noise $\varepsilon_i\sim N(0,1)$. Shapes are linear, concave (saturating), convex and sigmoid (threshold). We also use neighbour-count spillover, two-hop spillover, a degree-heterogeneous direct effect and
equilibrium propagation $Y=(I-\rho W)^{-1}(b+\tau z)+\varepsilon$ (Appendix~\ref{app:setup}).

\textbf{Estimators.}
The regression-based estimator fits OLS of outcomes on features $\phi_i(z)$ and predicts at the global treatment and control vectors:
$\hat\tau=\big(\tfrac1n\sum_i[\phi_i(\mathbf 1)-\phi_i(\mathbf 0)]\big)^{\!\top}\hat\beta$ \citep{chin2019,hanugander2023}.
Our main specification, \textbf{OLS$(z,f)$}, uses $\phi=(1,z_i,f_i)$; richer ones add the neighbour count, a two-hop fraction, log-degree
interactions, polynomials in $f$ (quad, cubic) or an 11-column set; a \textbf{correctly specified} (oracle) regression with the true
exposure columns is a control in every sweep, and gradient boosting is the flexible learner.
Baselines are the difference in means (DM) and Hájek exposure-mapping estimators (full-neighbourhood and $\kappa=0.75$) \citep{aronowsamii2017,chin2019}.

\textbf{Designs and graphs.}
Bernoulli$(p)$ assignment and graph-cluster randomisation with Louvain communities \citep{blondel2008} (Bernoulli$(0.5)$ per cluster), with random clusters of the same sizes and, on the SBM, planted blocks as controls.
Graphs: SBM \citep{holland1983}, Barab\'asi--Albert \citep{barabasi1999}, Watts--Strogatz \citep{watts1998} and Erd\H{o}s--R\'enyi (all $n=5000$, mean degree about 20; NetworkX \citep{hagberg2008}), and the real
Facebook100 Stanford3 \citep{traud2012,rossi2015} ($n=11{,}586$, mean degree 98) and SNAP ego-Facebook \citep{mcauley2012} ($n=4039$, mean degree 44); scaling grids vary $n\in[10^3,2\cdot10^4]$ and $d\in[5,80]$ (Appendix~\ref{app:setup}, \ref{app:repro}).

\textbf{Metrics.}
Per cell: bias, SD, RMSE, relative RMSE $=\mathrm{RMSE}/|\tau_G|$ and coverage of nominal 95\% Wald intervals (HC0, cluster-robust, graph-HAC over 1 or 2 hops), with
Monte-Carlo standard errors. The truth is exact, so $\mathrm{MSE}=\mathrm{bias}^2+\mathrm{SD}^2$ and the components are attributed separately.

\textbf{Where the error comes from.}
Three facts organise the study; the first two are made exact under stated conditions, the third is tested numerically.

\emph{(i) Bias is extrapolation of the local slope.} Under Bernoulli$(p)$ the own treatment $z_i$ is independent of $f_i$, which depends only on
neighbours. For noiseless one-hop outcomes the population least-squares coefficient on $f$ in OLS$(z,f)$ is
$s=\mathrm{Cov}(h(f),f)/\mathrm{Var}(f)$, the covariance being over the mixture of units with $f_i=\mathrm{Bin}(d_i,p)/d_i$. The estimator predicts $\tau+\gamma s$ whereas
the truth is $\tau+\gamma$, so
\begin{equation}\label{eq:bias}
  \mathrm{bias}_\infty \;=\; \gamma\,(s-1)\;\;\xrightarrow{\;d_i\to\infty\;}\;\;\gamma\,\big(h'(p)-1\big).
\end{equation}
The estimator transports the slope of the spillover curve at the operating point $p$ to the full range $[0,1]$. It is exact for $h$ linear (zero bias, for
any graph) and nonzero otherwise; it is a property of the model and the design and does not shrink with $n$ (Sec.~\ref{sec:res-theory}).
The derivation is in Appendix~\ref{app:theory}; for non-Bernoulli designs $z$ and $f$ are correlated and (\ref{eq:bias}) no longer applies.

\emph{(ii) Variance scales as $\sqrt{(1+d_h)/(n\,p(1-p))}$.} With errors of variance $\sigma_r^2$ treated as independent, $\mathrm{Var}(\hat\tau)\approx
\frac{\sigma_r^2}{n}\big[\frac1{p(1-p)}+\frac1{\mathrm{Var}(f)}\big]$ and $\mathrm{Var}(f)=p(1-p)\,\overline{d^{-1}}$, i.e.
$\mathrm{SD}\propto\sqrt{(1+d_h)/(n\,p(1-p))}$ with $d_h$ the harmonic mean degree. The exposure $f$ varies by only $O(1/\sqrt d)$ around $p$ whereas the
corners $f\in\{0,1\}$ lie $O(\sqrt d)$ standard deviations away, so the adjusted estimator pays a variance price that the difference in means does not.
The dependence among units that share neighbours changes the level, not the claimed scaling; both are tested.

\emph{(iii) Design and inference.} Cluster randomisation widens the support of $f$ (many units with $f\approx 0$ or $1$), which should reduce both
terms \emph{if} clusters follow the graph; and iid-style intervals ignore the dependence created by shared neighbours and by bias. We state the
predictions H1--H8 from the plan (Appendix~\ref{app:hyp}, which states when the numerical rules were fixed). Some of them (H1, H2, H3b, H6a, and the proportionality H3g of a later grid) are consequences of the model's algebra or checks of our code; we call them consistency checks and report them as passed or failed. The others are empirical, and are reported as supported, refuted or unclear.
